% ===== §2 =====
\section{The Epistemology Recalled}
\label{sec:recap}

\noindent
The account this paper applies can be recalled in its essentials, so that what follows stands on its own for a reader who has not the companion work to hand. The starting point is a limitation that traditional epistemology leaves unexamined, the supposition that a concept can be adequately carried by a single symbolic system, so that to know is to find the correct expression and the work of knowing is finished once that expression is secured. Against this supposition the epistemology holds that no symbolic system completes a concept. Every symbolization selects, compresses, and situates, and so renders a partial and situated aspect and leaves a remainder. A concept becomes observable through the relations among several such partial renderings, as a form is inferred from several projections none of which is the form.

From this the epistemology draws a term to stand in the place of representation. Knowing is generative observation. What is known is observed in the movement between symbolic expressions, the observation is an achievement that stands short of identity with the matter observed, and because the relations among expressions increase and change through history the observation is never final and the knowing continues. Three features of the account bear on what follows, and each is recalled here in turn.

The first is that observation is relational. The matter observed appears in no single rendering and comes into view only in the relations among several. The epistemic yield is located in the relations themselves, which are held to be productive of observation and not a mere connection of findings already secured. A relation, on this account, generates an observation that neither of its terms contained, and this generativity of relation is the deepest commitment of the account.

The second is that observation is generative in time as well as in structure. The field of relations through which a matter is observed is not a standing inventory awaiting discovery. It comes into being in history. New renderings arrive as the conditions of a form of life bring them forth, and each arrival opens relations that were unavailable before it, so that the observation of a matter develops as its field of relations grows. The account renders this historicity in the terms of a dialectical and materialist conception, on which the expressions and relations through which anything is observed arise from the conditions of a form of life and develop through the contradictions those conditions generate. Truth, so understood, is a becoming, generated in a historical and contradictory movement, and it holds none of the fixity of a possession.

The third is that observation is incomplete, and its incompleteness is not a defect a fuller state of knowledge would repair. It follows from the structure of the account, since the renderings of any matter are always finite in number and partial in reach, and there is always a further rendering, a further relation, a further aspect not yet brought into view. This incompleteness carries an epistemic humility that follows as a consequence of the account and forms no rhetorical modesty, and it will govern the treatment of understanding between persons, where the temptation to take one's observation of the other for the other's whole reality is the standing danger.

These three features, that observation is relational, that it is generative in time, and that it is incomplete, are the whole of what the following parts require. The reader who holds them can follow the application to intimacy without the companion work, and the reader who wishes the full argument will find it there.
